\documentclass[12pt,a4paper]{article}
\usepackage[utf8]{inputenc}
\usepackage{amsmath, amssymb}
\usepackage{geometry}
\usepackage{graphicx}
\usepackage{listings}
\usepackage{xcolor}
\usepackage{hyperref}

\geometry{top=2.5cm, bottom=2.5cm, left=2.5cm, right=2.5cm}

\title{\textbf{Selected Topics in Cryptography}\\ Lab Session 02: Toy ECDH}
\author{Estudiante}
\date{\today}

\lstset{
    language=C,
    basicstyle=\ttfamily\footnotesize,
    keywordstyle=\color{blue}\bfseries,
    commentstyle=\color{green!60!black},
    stringstyle=\color{red},
    showstringspaces=false,
    frame=single,
    numbers=left,
    numberstyle=\tiny,
    breaklines=true
}

\begin{document}

\maketitle

\section*{1. Personal Information}
\textbf{Course:} Selected Topics in Cryptography \\
\textbf{Date:} \today \\
\textbf{Topic:} Elliptic Curve Cryptography: Point Multiplication and ECDH Simulation.

\section*{2. Source Code}
The elliptic curve operations have been adapted to use \textbf{Standard Projective Coordinates} $(X, Y, Z)$ to represent affine points $(X/Z, Y/Z)$ avoiding costly modular inversions during addition and doubling operations.

\subsection*{Algorithm 1: Right-to-Left Point Multiplication}
\begin{lstlisting}
int elipar_point_mul_r2l(point_t* R, const point_t* P, const mpz_t k, const mpz_t a, const mpz_t b, const mpz_t p, bool print_steps) {
    point_t res, temp;
    point_init(&res); point_init(&temp);
    point_set_infinity(&res);
    point_set(&temp, P);
    
    mpz_t k_copy; mpz_init_set(k_copy, k);
    size_t i = 0;
    while (mpz_cmp_ui(k_copy, 0) > 0) {
        if (mpz_odd_p(k_copy)) {
            elipar_point_add(&res, &res, &temp, a, b, p);
        }
        elipar_point_double(&temp, &temp, a, b, p);
        mpz_fdiv_q_2exp(k_copy, k_copy, 1);
        i++;
    }
    
    point_set(R, &res);
    mpz_clear(k_copy); point_clear(&res); point_clear(&temp);
    return 0;
}
\end{lstlisting}

\subsection*{Algorithm 2: Left-to-Right Point Multiplication}
\begin{lstlisting}
int elipar_point_mul_l2r(point_t* R, const point_t* P, const mpz_t k, const mpz_t a, const mpz_t b, const mpz_t p, bool print_steps) {
    point_t res, temp_p;
    point_init(&res); point_init(&temp_p);
    point_set_infinity(&res);
    point_set(&temp_p, P);
    
    mpz_t k_copy; mpz_init_set(k_copy, k);
    size_t num_bits = mpz_sizeinbase(k_copy, 2);
    
    for (long i = num_bits - 1; i >= 0; i--) {
        elipar_point_double(&res, &res, a, b, p);
        if (mpz_tstbit(k_copy, i)) {
            elipar_point_add(&res, &res, &temp_p, a, b, p);
        }
    }
    
    point_set(R, &res);
    mpz_clear(k_copy); point_clear(&res); point_clear(&temp_p);
    return 0;
}
\end{lstlisting}

\section*{3. Testing Functions}
Below are 5 examples using both the Right-to-Left and Left-to-Right algorithms for $kP$. 
We used the curve $y^2 \equiv x^3 + 2x + 3 \pmod{97}$ with generator point $P = (3, 6)$. 
Notice that this specific point $P$ generates a subgroup of order 5 ($5P = \mathcal{O}$).

\begin{verbatim}
--- Test 1: k = 5 ---
>> Algoritmo R2L:
  [R2L i=0] k_i=1 -> Q = Q + P_i. Actual Q: (3, 6)
  [R2L i=1] k_i=0 -> Omitir suma.
  [R2L i=2] k_i=1 -> Q = Q + P_i. Actual Q: O
>> Algoritmo L2R:
  [L2R i=2] Q = 2Q. Actual Q: O
  [L2R i=2] k_i=1 -> Q = Q + P. Actual Q: (3, 6)
  [L2R i=1] Q = 2Q. Actual Q: (80, 10)
  [L2R i=1] k_i=0 -> Omitir suma.
  [L2R i=0] Q = 2Q. Actual Q: (3, 91)
  [L2R i=0] k_i=1 -> Q = Q + P. Actual Q: O
Resultado: O (Punto al infinito)

--- Test 2: k = 10 ---
>> Algoritmo R2L:
  [R2L i=0] k_i=0 -> Omitir suma.
  [R2L i=1] k_i=1 -> Q = Q + P_i. Actual Q: (80, 10)
  [R2L i=2] k_i=0 -> Omitir suma.
  [R2L i=3] k_i=1 -> Q = Q + P_i. Actual Q: O
Resultado: O

--- Test 3: k = 15 ---
>> Algoritmo R2L:
  [R2L i=0] k_i=1 -> Q = Q + P_i. Actual Q: (3, 6)
  [R2L i=1] k_i=1 -> Q = Q + P_i. Actual Q: (80, 87)
  [R2L i=2] k_i=1 -> Q = Q + P_i. Actual Q: (80, 10)
  [R2L i=3] k_i=1 -> Q = Q + P_i. Actual Q: O
Resultado: O

--- Test 4: k = 20 ---
>> Algoritmo R2L:
  [R2L i=0] k_i=0 -> Omitir suma.
  [R2L i=1] k_i=0 -> Omitir suma.
  [R2L i=2] k_i=1 -> Q = Q + P_i. Actual Q: (3, 91)
  [R2L i=3] k_i=0 -> Omitir suma.
  [R2L i=4] k_i=1 -> Q = Q + P_i. Actual Q: O
Resultado: O

--- Test 5: k = 25 ---
>> Algoritmo R2L:
  [R2L i=0] k_i=1 -> Q = Q + P_i. Actual Q: (3, 6)
  [R2L i=1] k_i=0 -> Omitir suma.
  [R2L i=2] k_i=0 -> Omitir suma.
  [R2L i=3] k_i=1 -> Q = Q + P_i. Actual Q: (3, 91)
  [R2L i=4] k_i=1 -> Q = Q + P_i. Actual Q: O
Resultado: O
\end{verbatim}

\section*{4. Simulation of ECDH}
An automated simulation was created for two entities (Alice and Bob) generating a shared secret $S$. We randomly generated a $128$-bit prime for the underlying finite field.

\begin{verbatim}
Parametros ECDH acordados:
p = 259724759763936159802421243210352256999
a = 192949947194419950194383621095450484007
b = 74149303083644888237766867364156607885
G = (5, 7)

Alice elige kA secreta: 224632427021277000312758535088198607810
Bob elige kB secreta: 100976876835531676369569759916640712836

Alice calcula y comparte A = kA * G:
  A = (173418783328457929096808771792846765109, 
       164068629066182190983474325829852729273)

Bob calcula y comparte B = kB * G:
  B = (50305030179271160272561851324741241275, 
       46913123138051945780540674811974877892)

Alice calcula el secreto S = kA * B:
  S = (60968394712849497951362774872976264714, 
       253273267623164229503681617562391842455)

Bob calcula el secreto S = kB * A:
  S = (60968394712849497951362774872976264714, 
       253273267623164229503681617562391842455)

=> ¡Las claves coinciden! El ECDH fue exitoso.
\end{verbatim}

\end{document}
